% Options for packages loaded elsewhere
\PassOptionsToPackage{unicode}{hyperref}
\PassOptionsToPackage{hyphens}{url}
\documentclass[
  ukrainian,
]{article}
\usepackage{xcolor}
\usepackage[margin=2.5cm]{geometry}
\usepackage{amsmath,amssymb}
\setcounter{secnumdepth}{-\maxdimen} % remove section numbering
\usepackage{iftex}
\ifPDFTeX
  \usepackage[T1]{fontenc}
  \usepackage[utf8]{inputenc}
  \usepackage{textcomp} % provide euro and other symbols
\else % if luatex or xetex
  \usepackage{unicode-math} % this also loads fontspec
  \defaultfontfeatures{Scale=MatchLowercase}
  \defaultfontfeatures[\rmfamily]{Ligatures=TeX,Scale=1}
\fi
\usepackage{lmodern}
\ifPDFTeX\else
\setmainfont{DejaVu Serif}
\setsansfont{DejaVu Sans}
\setmonofont{DejaVu Sans Mono}
\setmathfont{STIX Two Math}
\fi
\ifPDFTeX\else
  % xetex/luatex font selection
\fi
% Use upquote if available, for straight quotes in verbatim environments
\IfFileExists{upquote.sty}{\usepackage{upquote}}{}
\IfFileExists{microtype.sty}{% use microtype if available
  \usepackage[]{microtype}
  \UseMicrotypeSet[protrusion]{basicmath} % disable protrusion for tt fonts
}{}
\makeatletter
\@ifundefined{KOMAClassName}{% if non-KOMA class
  \IfFileExists{parskip.sty}{%
    \usepackage{parskip}
  }{% else
    \setlength{\parindent}{0pt}
    \setlength{\parskip}{6pt plus 2pt minus 1pt}}
}{% if KOMA class
  \KOMAoptions{parskip=half}}
\makeatother
\usepackage{graphicx}
\graphicspath{{/Users/kalinicenkovitalijmikolajovic/.codex/worktrees/5c15/Script_for_translate_book/Output_Final/}}
\makeatletter
\newsavebox\pandoc@box
\newcommand*\pandocbounded[1]{% scales image to fit in text height/width
  \sbox\pandoc@box{#1}%
  \Gscale@div\@tempa{\textheight}{\dimexpr\ht\pandoc@box+\dp\pandoc@box\relax}%
  \Gscale@div\@tempb{\linewidth}{\wd\pandoc@box}%
  \ifdim\@tempb\p@<\@tempa\p@\let\@tempa\@tempb\fi% select the smaller of both
  \ifdim\@tempa\p@<\p@\scalebox{\@tempa}{\usebox\pandoc@box}%
  \else\usebox{\pandoc@box}%
  \fi%
}
% Set default figure placement to htbp
\def\fps@figure{htbp}
\makeatother
\ifLuaTeX
\usepackage[bidi=basic,shorthands=off,provide=*]{babel}
\else
\usepackage[bidi=default,shorthands=off,provide=*]{babel}
\fi
\ifLuaTeX
  \usepackage{selnolig} % disable illegal ligatures
\fi
\setlength{\emergencystretch}{3em} % prevent overfull lines
\providecommand{\tightlist}{%
  \setlength{\itemsep}{0pt}\setlength{\parskip}{0pt}}
\usepackage{bookmark}
\IfFileExists{xurl.sty}{\usepackage{xurl}}{} % add URL line breaks if available
\urlstyle{same}
\hypersetup{
  pdflang={uk},
  hidelinks,
  pdfcreator={LaTeX via pandoc}}

\author{}
\date{}

\begin{document}

\section{Определение конформного
отображения}\label{ux43eux43fux440ux435ux434ux435ux43bux435ux43dux438ux435-ux43aux43eux43dux444ux43eux440ux43cux43dux43eux433ux43e-ux43eux442ux43eux431ux440ux430ux436ux435ux43dux438ux44f}

Отображение w=\(f(z)\) комплексиой плоскости z на комплексную плоскость
w называется конформным в точке \(z_{\bullet}\) плоскости z, если
производная \(f'(z_{\bullet}) \neq 0\) . Отображение \(f(z)\) называется
конформным в области D, если \(f'(z) \neq 0\) в каждой точке области D.

Например, отображение \(f(z) = z^2\) конформно всюду, за исключением
точки z = 0, поскольку \(z'(z) = 2z \neq 0\) для всех \(z \neq 0\) . C
другой стороны, отображение \(e^z\) конформно во всей плоскости z, так
как всюду \(f'(z) = e^z \neq 0\) . Какая же польза от конформных
отображений? Ответ состоит в том, что, решая уравнение Лапласа
\(\phi_{xx} + \phi_{yy} = 0\) внутри какой-то области плоскости
переменных x и y, мы можем рассматривать эту плоскость как плоскость
комплексного переменного z. Рассмотрим комплексное отображение w =
\(f(z)\) плоскости z на плоскость w, в которой введены координаты (u,
v). При отображении w = \(f(z)\) уравнение Лапласа
\(\varphi_{xx} + \varphi_{yy} = 0\) преобразуется в некоторое новое
уравнение с частными производными, зависящее от новых координат и и v.
Замечательным является следующее: если отображение w = \(f(z)\) является
конформным в области, где задано уравнение \(\phi_{xy} + \phi_{uy} = 0\)
, то новое уравнение также будет уравнением Лапласа в координатах и и v.
To есть уравнение \(\phi_{xx} + \phi_{uu} = 0\) перейдет снова в
уравнение \(\varphi_{nn} + \varphi_{nn} = 0\) . Так мы приходим к идее
найти такое конформное отображение, которое переводит область со сложной
границей в область с простой границей (напомним, что отображение
\(w=z^*\) переводит границу первого квадранта плоскости z в
действительную ось плоскости w).

Приведем несколько примеров, в которых будет показано, как можно решать
уравнения с частными производными с помощью конформных отображений.

Уравнение Лапласа в верхней полуплоскости.

Предположим, что нам необходимо решить следующую задачу Дирихле в
верхней полувлоскости (рис. 47.3):

(47.1)

\begin{verbatim}
можно найти стандартным методом разделения переменных. При этом оказывается, что решение   u\left(x,\,t\right)  , представленное в виде ряда
\end{verbatim}

Один із способів розв'язання цієї проблеми --- застосувати перетворення
Фур'є до змінної x, але набагато краще конформно відобразити верхню
половину площини у площину (u, v). Читачеві не буде складно побачити, що
конформне відображення

является суперпозицией простейших колебаний. Пространственная часть
\(X_n\)(x) каждого элементарного колебания является решением
(собственной функцией) некоторой краевой задачи Штурма---Лиувилля.

\begin{verbatim}
Итак, мы научились решать волиовое уравнение   u_{II}=c^2u_{xx}   на всей прямой   -\infty < x < \infty  . Формула Даламбера представляет решение в виде суммы двух бегущих волн, распространяющихся в противоположные стороны. Если же рассматривать это же уравнение на ограниченном промежутке 0 < x < L, то бегущих волн уже не станет, так как они будут взаимодействовать с границами. Вместо них возникают другие волны, которые называют стоячими волнами. Посмотрим, например, что случится, если гитарную струну, закрепленную в точках x=0, x=L, привести в движение. Для ответа на этот вопрос нам необходимо решить следующую смешанную задачу гиперболического типа:
\end{verbatim}

Вона має такі властивості (див. RNS. 47.4).

\begin{enumerate}
\def\labelenumi{\arabic{enumi}.}
\tightlist
\item
  Вона відображає верхню півплощину площини z у смуги
  \(-\infty < u < \infty\) , \(0 < v < \pi\) , у площині w.
\end{enumerate}

\pandocbounded{\includegraphics[keepaspectratio]{images/_page_357_Figure_11.jpeg}}

РИС 47.3 Задача граничного значення у верхній половині площини.

\begin{itemize}
\tightlist
\item
  \begin{enumerate}
  \def\labelenumi{\arabic{enumi}.}
  \setcounter{enumi}{1}
  \tightlist
  \item
    Ребро прямої y=0, -1 \textless{} x \textless{} 1, у площині z (де
    потенціал дорівнює 1), переходить у пряму \(v=\pi,\) -
    \(\infty < u < \infty\) , у площині w. - 3. Два промені
    x\textgreater1, y=0, і x\textless-1, y=0, у площині z проходять
    відповідно у позитивний і негативний
  \end{enumerate}
\end{itemize}

Читач повинен бути впевнений, що може підтвердити ці твердження.

\end{document}
